\section{Monodromy, restriction maps, and the component count}

\subsection{Abelianization and the corrected count}

\begin{proposition}[D.1.1; Proved]\label{st:D.1.1}
For an action $\rho:Q\to\operatorname{Aut}(K)$ and the corresponding
semidirect product $G=K\rtimes_\rho Q$, put
\[
(K^{\rm ab})_Q=K^{\rm ab}/
\langle\rho(q)\bar k-\bar k:q\in Q,\ \bar k\in K^{\rm ab}\rangle.
\]
There is an isomorphism
\[
G^{\rm ab}\longrightarrow Q^{\rm ab}\oplus(K^{\rm ab})_Q,
\qquad [(k,q)]\longmapsto([q],[k]).
\]
For a split extension without a specified splitting, the associated
short exact sequence is canonical and its splitting need not be.
\end{proposition}
\begin{proof}
The product is $(k,q)(k',q')=(k\rho(q)(k'),qq')$.
In the displayed quotient $[\rho(q)(k')]=[k']$, so the displayed map
is a surjective homomorphism to an abelian group. Its kernel contains
$[K,K]$, $[Q,Q]$ and all elements $\rho(q)(k)k^{-1}$.
These elements are themselves commutators in $G$. Conversely their
normal closure makes the images of $K$ and $Q$ abelian and makes
those images commute, so the quotient by that normal closure is
abelian. The kernel is exactly $[G,G]$. The inverse homomorphism
sends the two summands to the images of $(1,q)$ and $(k,1)$.
The induced maps in
$0\to(K^{\rm ab})_Q\to G^{\rm ab}\to Q^{\rm ab}\to0$
come from inclusion and projection; choosing the original group
section supplies its splitting.

For a concrete failure of the uncorrected direct sum, take
$K=Q=\Z$ and let the generator of $Q$ act on $K$ by $k\mapsto-k$.
Then the coinvariant quotient is $\Z/2\Z$, and the displayed map
gives $G^{\rm ab}=\Z\oplus\Z/2\Z$, whereas
$Q^{\rm ab}\oplus K^{\rm ab}=\Z^2$.
\end{proof}

\begin{proposition}[D.1.2; Proved]\label{st:D.1.2}
Let $f:Y\to C$ be a proper holomorphic surjection from a connected
compact complex threefold to a compact Riemann surface, with
connected fibres. Let $\Delta$ be a nonempty finite set containing
the critical values, $s=|\Delta|$, $g=g(C)$,
$D=f^{-1}(\Delta)_{\rm red}$, and let $r$ be the number of
irreducible components of $D$. Write $F$ for a smooth fibre and
\[
t'=\dim_{\Q}H_1(F;\Q)_{\pi_1(C\setminus\Delta)}
   =\dim_{\Q}H^1(F;\Q)^{\pi_1(C\setminus\Delta)}.
\]
For $b_2(Y)=0$ the exact component count is
\[
\boxed{\quad r=2g+s-1+t'-b_1(Y).\quad}
\]
Under CDP20's integral assumptions
$H^1(Y;\Z)=H^2(Y;\Z)=0$, the base is $\PP^1$ and
\[
r=s-1+t',\qquad
H_1(F;\Z)_{\pi_1(C\setminus\Delta)}\text{ is free abelian}.
\]
The last assertion concerns the indicated quotient of $H_1(F;\Z)$;
it does not assert that $H_1(F;\Z)$ itself is torsion-free.
\end{proposition}
\begin{proof}
Put $U=Y\setminus D$ and $C^\circ=C\setminus\Delta$.
The map $U\to C^\circ$ is a differentiable fibre bundle. Since
$\pi_2(C^\circ)=0$, its homotopy sequence is
\[
1\longrightarrow\pi_1(F)\longrightarrow\pi_1(U)
 \longrightarrow\pi_1(C^\circ)\longrightarrow1.
\]
The last group is free of rank $2g+s-1$. Lifting its free generators
gives a section of this extension. Proposition D.1.1 therefore gives
\[
H_1(U;\Z)\simeq
\Z^{2g+s-1}\oplus H_1(F;\Z)_{\pi_1(C^\circ)}.       \tag{D1}
\]
The second equality defining $t'$ follows by duality: the dual of
the coinvariant quotient is precisely the space of invariant
linear functionals on $H_1(F;\Q)$.

The cohomology sequence with compact supports for the open inclusion
$U\subset Y$ contains
\[
H^4(Y;\Q)\longrightarrow H^4(D;\Q)
\longrightarrow H_c^5(U;\Q)
\longrightarrow H^5(Y;\Q)\longrightarrow H^5(D;\Q)=0.
\]
The top-degree cohomology of the compact pure complex surface $D$
is $\Q^r$. One can see this by deleting the singular and intersection
loci, whose real dimensions are at most two, and taking compactly
supported top cohomology of its $r$ oriented smooth components.
Poincar\'e duality identifies $H_c^5(U;\Q)$ with $H_1(U;\Q)$ and
$H^5(Y;\Q)$ with $H_1(Y;\Q)$. Since $H^4(Y;\Q)=0$, taking
dimensions and using (D1) proves the formula.

For the integral refinement, the universal coefficient theorem shows
first that $H_1(Y;\Z)$ is finite, using $H^1(Y;\Z)=0$, and then that
it is zero, using
$\operatorname{Ext}^1(H_1(Y;\Z),\Z)\subset H^2(Y;\Z)=0$.
The group $H^4(Y;\Z)$ can be finite and nonzero. Its homomorphism
to $H^4(D;\Z)=\Z^r$ is nevertheless zero. The same sequence now
gives $H_1(U;\Z)\simeq\Z^r$, since $H^5(Y;\Z)=H_1(Y;\Z)=0$.
Equation (D1) then proves that the coinvariant summand is free.
Finally $H^1(C;\Q)\hookrightarrow H^1(Y;\Q)$ by the Leray edge
map, so $g=0$.
\end{proof}

\begin{example}[D.1.3; Exact check on the retained construction]
\label{st:D.1.3}
In the ordered homology basis
$(\widehat\gamma,\widehat u,\widehat w,\widehat\delta)$,
the two finite monodromies are
\[
A_1=\begin{pmatrix}1&0&0&0\\6&0&1&0\\-6&-1&-1&0\\-2&1&0&1\end{pmatrix},
\qquad
A_2=\begin{pmatrix}1&0&0&0\\0&0&-1&0\\-6&1&0&0\\3&0&1&1\end{pmatrix}.
\]
They satisfy $A_1^3=A_2^4=1$; the third generator is
$(A_1A_2)^{-1}$. The columns of
$[A_1-I\mid A_2-I]$ span the primitive subgroup
$\ker(\gamma:\Lambda\to\Z)$, where $\gamma=(1,0,0,0)$.
Indeed its first row is zero and its Smith diagonal is $(1,1,1,0)$.
The invariant homology vectors are $\Z\widehat\delta$, while the
invariant cohomology covectors are $\Z\gamma$. Consequently
\[
\Lambda_{\langle A_1,A_2\rangle}=\Z,
\qquad t'=1,
\qquad r=3-1+1=3.
\]
The printed count would give $r=3-1+4=6$.
These assertions are proved by exact integer matrix computation in
\texttt{code/D\_exact\_checks.py}; the Smith form and both invariant
spaces are printed in its output. The integral index in R7 is
verified next; the count uses the displayed quotient just computed.
\end{example}

\begin{proposition}[D.1.14; Proved]\label{st:D.1.14}
Put $e=\widehat\delta$,
$J_2=\left(\begin{smallmatrix}0&1\\-1&0\end{smallmatrix}\right)$,
and let $P_e$ be the integral stabilizer of $e$ for the form
\[
\Omega=\begin{pmatrix}0&0&0&6\\0&0&1&0\\0&-1&0&0\\-6&0&0&0\end{pmatrix}.
\]
Every element of $P_e$ has a unique expression
\[
g(B,x,t)=\begin{pmatrix}1&0&0\\6x&B&0\\t&-x^tJ_2B&1\end{pmatrix},
\quad B\in SL_2(\Z),\quad x\in\Z^2,\quad t\in\Z,
\]
in blocks of sizes $1,2,1$. Let
$Q(x)=x_1+x_2+x_1x_2$, and let
$\alpha:SL_2(\Z)\to\Z/12$ send
$S=\left(\begin{smallmatrix}0&-1\\1&0\end{smallmatrix}\right)$ to $9$
and $T=\left(\begin{smallmatrix}1&1\\0&1\end{smallmatrix}\right)$ to $1$.
Then
\[
\langle A_1,A_2\rangle
=\{g(B,x,t):t+6Q(x)-\alpha(B)=0\pmod{12}\},
\qquad [P_e:\langle A_1,A_2\rangle]=12.
\]
\end{proposition}
\begin{proof}
Preservation of $e$ and $\Omega$ forces the first row to be
$(1,0,0,0)$ and the last column to be $e$. The middle block has
determinant one. If its first column below the top entry is $v$,
the symplectic equations force the final middle row to be
$-v^tJ_2B/6$. Integrality and unimodularity of $B$ imply
$v\in6\Z^2$, proving the parametrization. Multiplication gives
\[
g(B,x,t)g(C,y,s)=g(BC,x+By,t+s-6x^tJ_2By).             \tag{D6}
\]
Set $a=\left(\begin{smallmatrix}0&1\\-1&-1\end{smallmatrix}\right)$.
Then $A_1=g(a,(1,-1)^t,-2)$,
$A_2=g(S,(0,-1)^t,3)$ and $Sa=T$.
The presentation
$SL_2(\Z)=\langle a,S:a^3=S^4=1,\ [a,S^2]=1\rangle$
follows from the Euclidean algorithm and the real projective-line
action. Modulo $\{\pm I\}$, $a$ has order three and $S$ order
two; $a,a^2$ send positive numbers into the two negative intervals
separated by $-1$, whereas $S$ sends negative numbers to positive
ones. A reduced alternating word cannot act trivially. Lifting
$S^2=-I$ gives the stated presentation. The values
$\alpha(a)=4$, $\alpha(S)=9$ respect the relations and give
$\alpha(T)=1$.

Modulo two, $Q$ is zero only at the origin. Therefore
$Q(Bx)=Q(x)$, and
$Q(x+y)-Q(x)-Q(y)=x^tJ_2y$ modulo two.
Equation (D6) makes
$\chi(g)=t+6Q(x)-\alpha(B)\pmod{12}$ a homomorphism.
It is onto because $\chi(g(I,0,1))=1$, and
$\chi(A_1)=\chi(A_2)=0$ by substitution.

For equality with its kernel, compute
\[
\begin{aligned}
K&=[A_1,A_2^2]=g(I,(0,-1)^t,6),\\
U&=A_2KA_2^{-1}=g(I,(1,0)^t,-6),\\
V&=K^{-1}=g(I,(0,1)^t,-6),\qquad Z=I+12e\gamma.
\end{aligned}
\]
One has $[U,V]=Z^{-1}$ and
\[
U^rV^sZ^n=g(I,(r,s)^t,-6(r+s+rs)+12n).
\]
Thus $U,V,Z$, all in the monodromy group, generate exactly the
intersection of $\ker\chi$ with the kernel of $g\mapsto B$.
The images $a,S$ generate $SL_2(\Z)$.
For an element of $\ker\chi$ choose a monodromy element with the
same middle block. Their quotient lies in that computed
intersection and hence in the monodromy group. This proves equality
and index twelve. The matrices and commutator are checked exactly
by the attached program.
\end{proof}

\subsection{The coherent restriction map and what its rank measures}

\begin{proposition}[D.1.4; Proved]\label{st:D.1.4}
For any $f:Y\to C$ as in D.1.2 put $E=R^1f_*\OO_Y$.
The coherent Leray sequence gives
\[
0\longrightarrow H^1(C,\OO_C)\xrightarrow{f^*}H^1(Y,\OO_Y)
\xrightarrow{e}H^0(C,E)\longrightarrow0.               \tag{D2}
\]
At a smooth fibre where cohomological base change holds, the
restriction homomorphism is exactly the composite
\[
H^1(Y,\OO_Y)\xrightarrow{e}H^0(C,E)
\xrightarrow{\operatorname{ev}_c}E\otimes\C(c)
\xrightarrow{\simeq}H^1(F_c,\OO_{F_c}).               \tag{D3}
\]
By Statement E4, $E$ is locally free. For $C=\PP^1$, write
$E=\bigoplus_{j=1}^q\OO(a_j)$.
Then the rank of (D3) at a general point is exactly
$\#\{j:a_j\geq0\}$. In particular, the nonnegative splitting
degrees claimed by CDP20 Proposition 5.5 are equivalent to its
claimed general-fibre surjectivity when $E$ has rank two and is
locally free.
\end{proposition}
\begin{proof}
Stein factorization gives $f_*\OO_Y=\OO_C$. A curve has no coherent
cohomology in degree two, so the low-degree Leray sequence is (D2).
The edge map is induced by restriction to inverse images of open
sets in $C$. Restricting these same representatives to $F_c$ gives
(D3), which proves the asserted identification of the actual maps.
Statement E4 proves local freeness by resolving the fibre support,
normalizing a cyclic base change with reduced central fibre, and
taking the invariant direct image. For a line bundle
$\OO(a)$ on $\PP^1$, its evaluation has rank one at every point
for $a\geq0$, and its space of global sections is zero for $a<0$.
Applying this separately to the direct summands gives the formula.
\end{proof}

\begin{proposition}[D.1.10; Proved]\label{st:D.1.10}
Let $f:Y\to\PP^1$ be a holomorphic algebraic reduction of a compact
complex threefold with $b_1(Y)=b_2(Y)=0$ and $a(Y)=1$.
At a general smooth fibre, the image of the restriction
$H^1(Y,\OO_Y)\to H^1(F,\OO_F)$ contains the complex span of
the coefficient homomorphism
\[
\iota_F:H^1(F;\Z)^{\pi_1(C^\circ)}
 \longrightarrow H^1(F,\OO_F),                         \tag{D4}
\]
which is induced by $\Z_F\hookrightarrow\OO_F$.
The final spanning step of CDP's rank argument reaches the whole
target exactly when this span is $H^1(F,\OO_F)$. Surjectivity of
the full restriction map need not have its image generated only by
this lattice, so that last equivalence describes the stated argument.
\end{proposition}
\begin{proof}
Let $D_1,\ldots,D_r$ be the prime components above $\Delta$, and
choose $N>0$ annihilating the finite group $H^2(Y;\Z)$.
The logarithms, in the exponential identification
$\Pic^0(Y)=H^1(Y,\OO_Y)$, of
$\OO_Y(ND_1),\ldots,\OO_Y(ND_r)$ generate a free abelian
subgroup of rank $r-s+1=t'$.
Here is the rank computation with its maps specified. The map is
$\Z^r\to\Pic^0(Y)$, $(n_i)\mapsto\OO_Y(N\sum n_iD_i)$.
A relation means that $N\sum n_iD_i$ is the divisor of a
meromorphic function. Since $f$ is the algebraic reduction, that
function is $f^*h$ for $h\in\mathcal M(\PP^1)$.
Its divisor is a sum of full fibres with total base degree zero.
Such relations have rank $s-1$, and those given by divisors of
rational functions supported on $\Delta$ already have this rank.
Statement D.1.2 gives the last equality.

Over $U=f^{-1}(C^\circ)$ each of these divisor line bundles has
its canonical trivialization. Exactness of the exponential sequence
therefore places its logarithm in the coefficient image of
$H^1(U;\Z)$. Coherent Leray identifies restriction
$H^1(Y,\OO_Y)\to H^1(U,\OO_U)$ with restriction of global
sections of $E$; this is injective because E4 makes $E$ locally free.
Thus the rank-$t'$ group of logarithms remains injective on $U$.

Topological Leray on the punctured curve gives
\[
0\to H^1(C^\circ;\Z)\to H^1(U;\Z)
\to H^1(F;\Z)^{\pi_1(C^\circ)}\to0.
\]
The first group's coefficient image vanishes because
$H^1(C^\circ,\OO)=0$. Conversely a class whose coefficient image
in $H^1(U,\OO)$ vanishes restricts to zero in $H^1(F,\OO_F)$;
Lemma D.1.7 makes its integral restriction to $F$ zero. Thus the
coefficient image of $H^1(U;\Z)$ identifies exactly with the
invariant lattice in (D4). The logarithm group maps to a subgroup
of the same rank $t'$, hence of finite index, in that lattice.
Its complex span equals the span of the entire lattice, and the
restriction image is a complex vector subspace. This proves the
assertion and the statement about the last spanning step.
\end{proof}

\begin{example}[D.1.5; Explicit second examples for the erroneous steps]
\label{st:D.1.5}
Let $E_0=F_0=E'=\C/(\Z+i\Z)$ and $a=1/2\in E_0$.
The involution
\[
(e,z)\longmapsto(e+a,-z)
\]
acts freely on $E_0\times F_0$. Put
$S=(E_0\times F_0)/\langle(a,-1)\rangle$,
$Y=S\times E'$, and $C=E_0/\langle a\rangle$.
The induced map $Y\to C$ is a smooth torus fibration with fibre
$F_0\times E'$ and monodromy
$\operatorname{diag}(-I_2,I_2)$ around the nontrivial deck character.
Its integral homology coinvariants are
$\Z^2\oplus(\Z/2\Z)^2$. After removing $s\geq1$ fibres,
Proposition D.1.1 gives
\[
H_1(U;\Z)=\Z^{s+3}\oplus(\Z/2\Z)^2,
\]
whereas the uncorrected direct-sum argument has free rank $s+5$.
This is an explicit compact complex realization of the faulty
group-theoretic step, independent of the retained six-sphere.

Let $\eta$ be the nontrivial order-two degree-zero line bundle on
$C$ associated to the double cover $E_0\to C$. Finite-cover descent
of fibre cohomology gives
\[
R^1f_*\OO_Y=\eta\oplus\OO_C,\qquad R^2f_*\OO_Y=\eta.
\]
The invariant part of the exterior algebra on the three
antiholomorphic one-form directions, with signs $(1,-1,1)$,
gives $h^q(Y,\OO_Y)=(1,2,1,0)$. The restriction from
$H^1(Y,\OO_Y)$ to the two-dimensional $H^1(F_0\times E',\OO)$
has image precisely the $E'$ direction and has rank one.
All these homomorphisms are induced by the displayed finite
\'etale quotient and projection. This example has $b_1(Y)=4$;
it is a counterexample to the unchecked restriction argument, not
a second example satisfying the integral CDP20 hypotheses.

There is also an explicit negative summand over $\PP^1$. Let
$h:R\to\PP^1$ be the relatively minimal elliptic surface with
Weierstrass equation
\[
y^2=x^3+(t^4+1)x+t^6.
\]
Its fundamental line bundle is $\OO(1)$ and its discriminant has
twelve simple zeroes. Relative duality gives
$R^1h_*\OO_R=\OO(-1)$. For $Y'=R\times E'\to\PP^1$,
the product K\"unneth maps give
\[
R^1f_*\OO_{Y'}=\OO\oplus\OO(-1),\qquad
R^2f_*\OO_{Y'}=\OO(-1),
\]
and the restriction in degree one again has rank one. This
projective example has $b_1=2$, $b_2=11$ and $a=3$, so it does not
replace the missing second counterexample under all the standing
CDP20 hypotheses.
\end{example}

\begin{example}[D.1.6; A second count counterexample under the integral hypotheses]
\label{st:D.1.6}
Wang's attached construction $X_2(\lambda)$ has integral homology
$H_0=H_6=\Z$, $H_2=H_3=\Z/2\Z$, and the other reduced
homology groups zero. Hence its integral cohomology satisfies
$H^1=H^2=0$. Its fibration has four exceptional values, over which
the reduced fibres are irreducible: the multiplicity-three
bielliptic fibre, two rank-one fillings, and the rank-two filling.
In the rank-one filling over the $I_2$ value, the primitive period
exponent is $(2,1)$; the component count is one, not two.
These are the construction's statements
\texttt{intro:main}, \texttt{fill:rankone},
\texttt{fill:a2}, \texttt{fill:log}, and
\texttt{main:theorem} in the attached primary source
\cite{Wang2026}.

In its cohomology basis the four displayed monodromies are
\[
\begin{aligned}
B_*&=\begin{pmatrix}1&0&-6&-2\\0&-1&-1&1\\0&1&0&-1\\0&0&0&1\end{pmatrix},&\qquad
B_2&=\begin{pmatrix}1&0&6&3\\0&1&2&1\\0&0&1&0\\0&0&0&1\end{pmatrix},\\
B_+&=\begin{pmatrix}1&0&0&0\\0&2&1&0\\0&-1&0&0\\0&0&0&1\end{pmatrix},&\qquad
B_-&=\begin{pmatrix}1&0&0&-1\\0&1&1&0\\0&0&1&0\\0&0&0&1\end{pmatrix}.
\end{aligned}
\]
Their product is one and their common invariant cohomology group
has rank one. On homology the matrices are $B_i^{-t}$; the Smith
diagonal of their difference matrix is $(1,1,1,0)$. Consequently
$t'=1$ and the corrected count is $r=4-1+1=4$, while the printed
count gives seven. The exact matrix calculation is included in the
program. The topological construction and its local fillings are
being used here as cited primary-source results, not inferred from
the matrices alone. Its coherent direct images are derived separately
in Statement D.1.13 below.
\end{example}

\subsection{A simultaneous non-torsion class on every full fibre}

\begin{lemma}[D.1.7; Proved]\label{st:D.1.7}
For a connected compact complex space $S$, including a nonreduced
one, the coefficient map
\[
H^1(S;\mathbb R)\longrightarrow H^1(S,\OO_S)
\]
is injective as a map of real vector spaces.
\end{lemma}
\begin{proof}
First take $S$ reduced. Represent a real cohomology class by a
locally constant real \v Cech cocycle $a_{ij}$. If its image is zero,
after refining the open cover there are holomorphic functions $u_i$
with $a_{ij}=u_i-u_j$. The imaginary parts of the $u_i$ agree on
overlaps and give a global continuous pluriharmonic function.
Pull it back to a compact resolution of each irreducible component.
The maximum principle makes each such pullback constant; intersection
of components and connectedness of $S$ make the constants equal.
Thus, after subtracting one imaginary constant from all $u_i$,
each $u_i$ is real-valued and holomorphic, hence locally constant.
They form a real coboundary for $a_{ij}$. For nonreduced $S$, a
class zero in $H^1(S,\OO_S)$ also becomes zero after applying
$\OO_S\to\OO_{S_{\rm red}}$. The underlying topological space
and its real cohomology do not change under reduction, so the
reduced assertion proves the result.
\end{proof}

\begin{lemma}[D.1.8; Proved]\label{st:D.1.8}
Let $f:Y\to C$ be as in D.1.2, and let $j:C^\circ\hookrightarrow C$
be the smooth-fibre locus. The natural generization homomorphism
\[
R^1f_*\Z_Y\longrightarrow
j_*\bigl((R^1f_*\Z_Y)|_{C^\circ}\bigr)
\]
is injective. Every nonzero global section of $R^1f_*\Z_Y$ is
nonzero in $H^1(f^{-1}(c);\Z)$ at every $c\in C$.
\end{lemma}
\begin{proof}
It suffices to check a critical value $c$. Represent a stalk class
by $a\in H^1(V;\Z)$ for $V=f^{-1}(\Delta_c)$ and a sufficiently
small disc $\Delta_c$. Put $V^*=f^{-1}(\Delta_c\setminus\{c\})$.
Assume that the restriction of $a$ to a smooth fibre is zero.
The homotopy sequence for $V^*\to\Delta_c^*$ shows that
$a|_{V^*}$ is the pullback of an integer multiple $k$ of the
positive meridian class of the punctured disc.

Choose a smooth point of one irreducible component of the central
divisor, away from all its other components. There the map has
the form $t=z^m u$ with $m>0$ and $u$ a holomorphic unit.
A small positive normal circle has base winding $m$ and bounds
its normal disc in $V$. Evaluating $a$ on that circle gives
$km=0$, whence $k=0$.
The inclusion $V^*\hookrightarrow V$ induces a surjection of
fundamental groups: a loop can be perturbed off a complex divisor,
including its higher-codimension singular strata. Therefore
$a|_{V^*}=0$ implies $a=0$.
Taking the direct limit over the discs proves the asserted stalk
injection; proper topological base change identifies its source
with $H^1(f^{-1}(c);\Z)$.

If a global section vanished at one point, it would vanish on a
neighbourhood of that point. Its restriction to the connected
smooth locus is a section of a local system, so it would vanish
everywhere on that locus. The just-proved injection then makes
it zero everywhere, as required.
\end{proof}

\begin{proposition}[D.1.9; Proved]\label{st:D.1.9}
Let $Y$ be a connected compact complex manifold with
$b_1(Y)=b_2(Y)=0$, and let $f:Y\to\PP^1$ be a proper
holomorphic surjection with connected fibres. There is a class
$\beta\in H^1(Y,\OO_Y)$ such that, for every irrational real number
$r$, the line bundle
\[
L_r=\exp(2\pi i\,r\beta)
\]
has infinite order on every scheme-theoretic fibre of $f$.
Every connected component of $\Pic(Y)$ likewise contains a real
one-parameter family with this property at every irrational real
parameter.
No local-freeness assumption on $R^1f_*\OO_Y$ or constant-dimension
assumption on fibre cohomology is required.
\end{proposition}
\begin{proof}
The low-degree integral Leray sequence and connectedness of the
fibres give an injection with finite cokernel
\[
d_2:H^0(\PP^1,R^1f_*\Z_Y)
 \hookrightarrow H^2(\PP^1;\Z)=\Z.                   \tag{D5}
\]
Indeed $H^1(Y;\Z)=0$ and $H^2(Y;\Z)$ is finite.
Thus the image in (D5) is $d\Z$ for an integer $d>0$; for the
stronger integral hypothesis $H^2(Y;\Z)=0$ one has $d=1$.
Choose $\gamma$ mapped to $d$. By Lemma D.1.8,
each $\gamma_c\in H^1(f^{-1}(c);\Z)$ is nonzero.
The coefficient morphism $\Z_Y\to\OO_Y$ sends $\gamma$ to a
global section of $R^1f_*\OO_Y$. Coherent Leray and
$H^1(\PP^1,\OO)=H^2(\PP^1,\OO)=0$ identify this space with
$H^1(Y,\OO_Y)$; let $\beta$ be its inverse image.

For every full fibre $Y_c$, naturality of the edge maps and of
restriction gives
\[
\beta|_{Y_c}=\iota_c(\gamma_c)\in H^1(Y_c,\OO_{Y_c}), 
\]
where $\iota_c$ is the coefficient map. This uses the natural
coherent base-change homomorphism, but does not require it to be
an isomorphism. If $L_r|_{Y_c}$ had finite order, the exponential
sequence on the full fibre would give an integer $n>0$ and
$\lambda\in H^1(Y_c;\Z)$ with
\[
nr\,\iota_c(\gamma_c)=\iota_c(\lambda).
\]
Lemma D.1.7 implies $nr\gamma_c=\lambda$ in real cohomology.
The group $H^1(Y_c;\Z)$ is a finitely generated free abelian group,
and $\gamma_c$ is a nonzero element of it. A nonzero integral
coordinate of $\gamma_c$ now implies that $r$ is rational, a
contradiction. This works for every $c$ and proves the simultaneous
assertion in $\Pic^0(Y)$.

For a component represented by a line bundle $A$, choose $N>0$
with $A^N\in\Pic^0(Y)$ and write $A^N=\exp(2\pi i b)$.
Then
\[
A_r=A\otimes\exp\bigl(2\pi i(r\beta-b)/N\bigr)
\]
belongs to that component and satisfies $A_r^N=L_r$. A group
element has finite order exactly when any fixed positive power of
it has finite order. Therefore each $A_r|_{Y_c}$ is non-torsion
for every irrational real $r$, proving the componentwise statement.
\end{proof}

\begin{proposition}[D.1.11; Proved]\label{st:D.1.11}
Let $f:Y\to\PP^1$ be a proper connected-fibre holomorphic map from
a connected compact complex manifold with $b_1(Y)=b_2(Y)=0$.
In every component of $\Pic(Y)$, the line
bundles which are non-torsion on every full fibre form a dense
$G_\delta$ set of full Lebesgue measure in its complex vector-space
coordinates. If evaluation of global sections of $R^1f_*\OO_Y$
has rank one, then $\Pic^0(Y)=\C$ and the exceptional set is
countable. For two-torus fibres this alternative holds whenever
the splitting has a negative summand.
In particular the general-line-bundle conclusion of CDP20
Proposition 5.6 holds with this precise quantifier, without the
nonnegative-summand assertion of Proposition 5.5.
\end{proposition}
\begin{proof}
Let $E=R^1f_*\OO_Y$, a vector bundle by E4, and use
$\gamma,\beta$ from D.1.9. The section of $E$ corresponding to
$\beta$ has no zero. In fact its value maps naturally to the
nonzero class $\iota_c(\gamma_c)$ in the cohomology of every full
fibre; Lemmas D.1.7 and D.1.8 prove nonvanishing. Write
$E=\bigoplus_j\OO(a_j)$ and let $k$ count the nonnegative degrees.
Evaluation has constant rank $k$ onto their direct-summand subbundle
$E_+$. Nonvanishing of $\beta$ gives $k\geq1$. When $k=1$, its
sole nonnegative degree is zero, since a positive-degree line bundle
has no nowhere-zero section. In that case $H^0(E)=\C\beta$.

Suppose first $k\geq2$, and put $V=H^1(Y,\OO_Y)=H^0(E)$ and
$N=\dim_{\C}V$. On a simply connected disc in the smooth-fibre
locus where degree-one base change holds, trivialize the integral
local system. For a flat integral
class $\lambda$, the equation
\[
z|_{F_c}=\iota_c(\lambda),\qquad (z,c)\in V\times\Delta
\]
defines an analytic set. The parameter values where
$\iota_c(\lambda)\in(E_+)_c$ form an analytic subset of the disc;
over it the fibres are affine spaces of dimension $N-k$.
The incidence set thus has complex dimension at most
$N-k+1\leq N-1$. Its projection to $V$ has Lebesgue
measure zero. It is also meagre: stratify and exhaust that set
by countably many compact sets; the images are closed sets with
empty interior, by their real dimension at most $2N-2$.
Take countably many discs, integral classes and positive integer
denominators to include every torsion condition.

There are finitely many excluded fibres, including the exceptional
values of degree-one base change. At each, the natural map
$E\otimes\C(c)\to H^1(Y_c,\OO_{Y_c})$ is injective. This follows
directly from the long exact sequence for multiplication by a local
base parameter: its cokernel is the parameter-torsion in
$R^2f_*\OO_Y$. The restriction $V\to H^1(Y_c,\OO_{Y_c})$
therefore has rank $k\geq2$. The torsion conditions are contained in
countably many affine subspaces of complex codimension at least
two. Combining these with the preceding projections gives a
meagre null set. The complement is dense and $G_\delta$, as claimed.

It remains to prove the case $k=1$.
Now $V=\C\beta$. On a simply connected disc in the same open locus,
the projection of a flat integral class $\lambda$ to
$E/(\OO\beta)$ is holomorphic. Unless it vanishes identically,
its zero set is discrete. Such zeroes give only countably many
parameters $z$ for which $\iota_c(\lambda)=z\beta_c$.
We prove that an identically vanishing projection gives a constant
parameter, including after analytic continuation.

Pass to the universal cover of the smooth locus, on which
$H^1(F;\Z)$ is marked. The identity theorem continues the relation
$\iota_c(\lambda)\in\C\iota_c(\gamma)$ everywhere on this cover.
If $\lambda$ and $\gamma$ are real-dependent, their ratio is a
constant rational number. Otherwise let
$W_{\Q}=\Q\gamma+\Q\lambda$. The coefficient map
$H^1(F;\R)\to H^1(F,\OO_F)$ is real-linearly injective by D.1.7.
The images of the two real-independent classes span the entire
real two-plane underlying $\C\beta_c$, so $W_{\R}$ is precisely
its inverse image. Analytic continuation around a base loop
sends $\lambda$ to a class whose coefficient image again belongs
to this line. Hence it preserves $W_{\Q}$ and fixes $\gamma$.

Choose an integral basis $(\gamma_0,\lambda_0)$ of
$W_{\Q}\cap H^1(F;\Z)$ with $\gamma$ a positive integral multiple
of $\gamma_0$. The induced monodromy on this lattice fixes
$\gamma_0$ and preserves the orientation of $W_{\R}$: the complex
structure induced on that plane by the coefficient map varies
continuously along the loop, so its complex orientation cannot change.
It therefore sends $\lambda_0$ to $\lambda_0+k\gamma_0$ for an
integer $k$. Put
\[
\tau(c)=\frac{\iota_c(\lambda_0)}{\iota_c(\gamma_0)}.
\]
This holomorphic function on the universal cover never takes a
real value, by real injectivity of the coefficient map, and its
imaginary part has constant sign. Monodromy changes $\tau$ by
integers. Accordingly either $\exp(2\pi i\tau)$ or
$\exp(-2\pi i\tau)$ descends to a holomorphic function of absolute
value less than one on $\PP^1$ minus finitely many points. The
removable singularity theorem extends it to $\PP^1$, where it is
constant. Its logarithmic derivative gives $d\tau=0$. Thus the
original ratio is constant as asserted.

There are countably many integral classes on a countable disc
cover. The isolated-zero alternatives and the constant alternatives
just proved give a countable exceptional set on smooth fibres.
Each exceptional full fibre contributes another countable set,
because $z\mapsto z\beta_c$ is injective and its intersection with
the rational integral coefficient image is countable. This proves
the assertion in $\Pic^0(Y)$.

Finally for any component $A\Pic^0(Y)$, the map taking the $N$th
power, for an integer $N$ killing its component class, is the affine
isomorphism $z\mapsto b+Nz$ of the complex parameter space used
in D.1.9. It preserves the torsion condition in every fibre, and
preserves countability, meagreness and measure-zero sets. This
proves all the componentwise assertions.
\end{proof}

The quantifier in D.1.11 is explicit. In the globally generated
two-torus case, the torsion parameters at even one fixed smooth fibre are
dense: the rational period lattice is dense in
$H^1(F,\OO_F)$, and restriction onto that space is surjective.
Thus the good set in that case cannot contain a nonempty open
subset of the Picard parameter space. A dense $G_\delta$ set of
full measure is the appropriate conclusion of the incidence
argument, rather than a Zariski-open conclusion.

\begin{proposition}[D.1.12; Proved]\label{st:D.1.12}
For a proper connected-fibre holomorphic map $f:Y\to C$ from a
compact complex threefold with $b_2(Y)=0$ and general fibre a
two-torus or a primary Kodaira surface, the cup-product homomorphism
is an isomorphism
\[
\bigwedge^2R^1f_*\OO_Y\xrightarrow{\simeq}R^2f_*\OO_Y.
\]
Every full fibre has structure-sheaf cohomology dimensions
$(h^0,h^1,h^2)=(1,2,1)$, and these two direct images satisfy
base change at every fibre. This proves these two cases of CDP20
Proposition 5.1 by an index and cup-product argument.
\end{proposition}
\begin{proof}
Write $E=R^1f_*\OO_Y$, which is locally free of rank two by E4,
and $G=R^2f_*\OO_Y$, which has rank one. Put
$d=\deg E$, $g=\deg(G/G_{\rm tor})$, and
$\ell=\operatorname{length}(G_{\rm tor})$.
Since $c_1(Y)=0$ rationally, Hirzebruch--Riemann--Roch gives
$\chi(Y,\OO_Y)=\int_Yc_1c_2/24=0$.
Coherent Leray and Riemann--Roch on $C$ give, with the rank terms
cancelled,
\[
0=\chi(\OO_C)-\chi(E)+\chi(G)=g+\ell-d.
\]
The cup product $\det E\to G$ is an isomorphism on a general
fibre. For a two-torus this is its exterior-algebra description.
For a primary Kodaira surface, use invariant $(1,0)$-forms
$\alpha,\eta$ with $d\alpha=0$ and
$d\eta=c\alpha\wedge\overline\alpha$, $c\ne0$.
Their conjugates are $\overline\partial$-closed. The class
$[\overline\alpha\wedge\overline\eta]$ is nonzero because its
Serre pairing with the holomorphic two-form $\alpha\wedge\eta$
is the nonzero integral of a volume form. This also proves that
$[\overline\alpha],[\overline\eta]$ are independent; since
$h^1(\OO)=2$ and $h^2(\OO)=1$, it proves the required cup-product
isomorphism. The invariant forms and their differential are
the forms in \cite[\S3, equations (9)--(12)]{Borcea1984}.
The composite of $\det E\to G$ with the torsion-free quotient is a
nonzero homomorphism of line bundles, so $g\geq d$. The displayed
identity and $\ell\geq0$ force $g=d$, $\ell=0$. A nonzero map
of line bundles of equal degree has no zero divisor, proving the
claimed global isomorphism.

Top-degree base change gives
$G\otimes\C(c)\simeq H^2(Y_c,\OO_{Y_c})$, of dimension one.
The multiplication-parameter exact sequence then identifies
$E\otimes\C(c)$ with $H^1(Y_c,\OO_{Y_c})$, since $G$ has no
parameter-torsion. E4 gives $h^0=1$. These arguments retain the
scheme-theoretic fibres throughout.
\end{proof}

\begin{proposition}[D.1.13; Refuted printed surjectivity and positivity, with a second example]
\label{st:D.1.13}
For Wang's threefold $Y=X_2(\lambda)$ in D.1.6 one has
\[
K_Y\simeq f^*\OO_{\PP^1}(-1),\qquad
R^1f_*\OO_Y\simeq\OO\oplus\OO(-1),\qquad
R^2f_*\OO_Y\simeq\OO(-1).
\]
Consequently $h^q(Y,\OO_Y)=(1,1,0,0)$, and restriction to every
smooth two-torus fibre has rank one. This threefold satisfies
$H^1(Y;\Z)=H^2(Y;\Z)=0$, $a(Y)=1$, and is a second
counterexample to CDP20 Corollary 5.2 and Proposition 5.5.
Its integral homology torsion distinguishes it from the retained
integral homology sphere.
\end{proposition}
\begin{proof}
We compute the canonical line through the actual local fillings.
The rational elliptic surface $\pi:S\to\PP^1$ in this construction
has Weierstrass equation
\[
y^2=x^3-3(1+t)x+(2+3t+2t^2)
\]
and fundamental line bundle $\OO(1)$; hence
$K_S=\pi^*\OO(-1)$. On the complement of the zero section in
$M=\OO_S(P-O)$, the determinant of the relative cotangent sequence
identifies its canonical line with the pullback of $K_S$.
In a local multiplicative coordinate $v$, this identification sends
a local two-form $\omega_S$ to $\omega_S\wedge d\log v$.
Changing $v$ by a nonzero holomorphic function on $S$ adds a
one-form from $S$ to $d\log v$, whose wedge with $\omega_S$ is zero.
This verifies the coordinate independence of the identification.
The contracting generator covers translation by $-6P$ on the
elliptic fibres. Translation preserves the relative elliptic
differential, and the extra base terms disappear after wedging
with the base differential. Its line-bundle lift changes $d\log v$
only by a base one-form. It therefore preserves this canonical
identification. The latter descends to
$K_Y=f^*\OO(-1)$ on the open quotient.

At either rank-one filling, let $q$ be the base coordinate and
$(x,y)$ the two multiplicative coordinates of the construction.
At $t=0$ the explicit change is $x=v$, $y=v^2/z$; its exponent
matrix is unimodular. At $\alpha_+$ no change is required.
Thus the preceding three-form is a nonzero holomorphic unit times
\[
dq\wedge d\log x\wedge d\log y.
\]
The nondegenerating $y$ direction becomes an elliptic curve after
quotienting by its contracting period. Its logarithmic differential
is a nonvanishing relative one-form; changes by a base differential
do not affect the displayed wedge. At a double curve the equation
is $uv=q$, and
$dq\wedge d\log u=-du\wedge dv$. Therefore the three-form extends
without a zero or pole over the full rank-one filling.

At $\alpha_-$ the rank-two filling uses a unimodular period lattice
and the standard toric charts $q=z_0z_1z_2$. A pair of independent
fibre characters, together with $q$, is a unimodular system of
characters in each smooth chart. The same form is, up to a nonzero
constant and a holomorphic unit,
\[
q\,d\log z_0\wedge d\log z_1\wedge d\log z_2
=dz_0\wedge dz_1\wedge dz_2.
\]
It extends without zero or pole, including at the triple points.

It remains to check the order-three logarithmic filling at infinity.
Put $s=1/t$, $x_s=s^2x$ and $y_s=s^3y$. Its Weierstrass
coefficients have orders three and four in $s$, giving type $IV^*$.
On the cubic base change $s=r^3$, put
$x_s=r^4X$, $y_s=r^6Y$. The good-reduction equation is
\[
Y^2=X^3-3r(r^3+1)X+(2+3r^3+2r^6),
\]
whose central elliptic curve is smooth. The canonical form satisfies
\[
ds\wedge\frac{dx_s}{y_s}
=3\,dr\wedge\frac{dX}{Y}.
\]
It has no zero or pole. For $\zeta=e^{2\pi i/3}$ the root deck
action has weights $r\mapsto\zeta r$, $X\mapsto\zeta^2X$,
$Y\mapsto Y$; the weights on $dr$ and $dX/Y$ multiply to one.
The additional invariant translation used in the logarithmic
transform preserves the holomorphic two-form on the torus cover.
Hence the three-form is invariant and descends through the free
cyclic action as a nonvanishing three-form. This proves that the
canonical identification from the open quotient extends everywhere,
giving $K_Y=f^*\OO(-1)$.

The map is flat and Gorenstein, since its source is smooth and its
base is a smooth curve. Relative Serre duality gives
\[
(R^2f_*\OO_Y)^*=f_*\omega_{Y/\PP^1}
=f_*f^*\OO(1)=\OO(1).
\]
Statement D.1.12 gives local freeness and
$\det(R^1f_*\OO_Y)=R^2f_*\OO_Y=\OO(-1)$.
The nowhere-zero section from D.1.9 embeds $\OO$ as a subbundle
of $R^1f_*\OO_Y$ with quotient $\OO(-1)$. The extension splits
because $H^1(\PP^1,\OO(1))=0$. Coherent Leray now computes all
four displayed structure-sheaf cohomology groups, and D.1.4
identifies the degree-one restriction with rank-one evaluation.

The homology computation in D.1.6 and the universal coefficient
theorem give $H^1(Y;\Z)=H^2(Y;\Z)=0$; specifically the torsion
in $H_2$ contributes to $H^3$, not to $H^2$.
The general fibre is K\"ahler and $b_2(Y)=0$, so E1 proves that
this map is the algebraic reduction and $a(Y)=1$.
The retained six-sphere has zero second integral homology, whereas
$H_2(Y;\Z)=\Z/2\Z$. An isomorphism of complex manifolds would
induce an isomorphism of these groups, which is impossible.
This completes the independent counterexample.
\end{proof}

\begin{example}[D.1.15; Proved comparison with identical integral monodromy]
\label{st:D.1.15}
Two smooth proper families of two-tori over $\PP^1$ have the same
constant marked $H^1(-;\Z)$ local system but have
$R^1f_*\OO=\OO^2$ and $R^1f_*\OO=\OO(1)^2$, respectively.
Thus the integral representation alone does not determine the
coherent splitting for general compact torus fibrations. Both
examples have $b_1=4$; this comparison does not assert a new
counterexample under $b_1=b_2=0$.
\end{example}
\begin{proof}
Let $V=\OO_{\PP^1}(1)^{\oplus2}$. On the affine coordinate $t$,
four global sections of $V$ are the columns of
\[
P(t)=\begin{pmatrix}1&i&t&it\\-t&it&1&-i\end{pmatrix}.
\]
They are a real basis in every fibre. Indeed, writing
$z_1=a+ib$, $z_2=c+id$, their real linear combination is
$(z_1+t z_2,\overline z_2-t\overline z_1)$.
Its vanishing gives $(1+|t|^2)\overline z_2=0$ and then $z_1=0$.
At infinity, the $\OO(1)$ frame changes $P(t)$ to $t^{-1}P(t)$,
whose leading columns again form a real basis. The lattice
generated by these sections therefore acts freely and properly by
holomorphic fibre translations on $V$. Its quotient $Z$ is a
compact smooth holomorphic torus bundle over $\PP^1$.
The displayed real basis identifies it differentiably with
$\PP^1\times(\R^4/\Z^4)$ and marks its constant integral local
system.

The fibrewise period map includes the bundle of holomorphic
one-forms by $P(t)^t$. The Dolbeault coefficient map gives the
exact sequence
\[
0\longrightarrow\OO(-1)^2
 \xrightarrow{P(t)^t}\OO^4
 \longrightarrow R^1f_*\OO_Z\longrightarrow0.
\]
The last map can be represented explicitly by
\[
B(t)=\begin{pmatrix}-it&-t&i&1\\-i&1&-it&t\end{pmatrix}.
\]
The identity $B(t)P(t)^t=0$ holds polynomially. Its two rows have
rank two at every finite $t$, as is seen from their two-by-two
minors, and their leading coefficients have rank two at infinity.
Thus this is the exact complex of vector bundles
\[
0\longrightarrow\OO(-1)^2\longrightarrow\OO^4
 \xrightarrow{B}\OO(1)^2\longrightarrow0,
\]
proving $R^1f_*\OO_Z\simeq\OO(1)^2$. The polynomial identity is
also checked in the attached exact program.

For the constant family $Z_0=Z_{t=0}\times\PP^1$, the same
marked fibre lattice gives the identical constant integral local
system, whereas coherent K\"unneth gives
$R^1f_*\OO_{Z_0}=H^1(Z_{t=0},\OO)\otimes\OO=\OO^2$.
The different splittings are therefore compared by an explicit
identity of the integral representations and explicit coherent
quotient maps.
\end{proof}

\subsection{Trivial monodromy (D.2)}

\begin{proposition}[D.2.1; Proved obstruction]\label{st:D.2.1}
A compact complex threefold with $b_1=b_2=0$ and a holomorphic
connected-fibre map to a curve, with general fibre a two-torus and
trivial rational monodromy, has reducible fibres. More precisely,
for every nonempty set of $s$ removed values containing all critical
values, their reduced fibres have $r=s+3$ irreducible components in
total. Every component of each reducible fibre has no compact
K\"ahler resolution. Consequently no such fibration can have all
its fibre components admitting compact K\"ahler resolutions.
\end{proposition}
\begin{proof}
The base is $\PP^1$ by D.1.2. Trivial monodromy on the
four-dimensional $H_1(F;\Q)$ gives $t'=4$, so the corrected formula
gives $r=s+3$. If all removed fibres were irreducible one would have
$r=s$, which is impossible. The assertion about each component is
Statement C.3.4, applied to this threefold. The last assertion follows
by combining these two conclusions.
\end{proof}

\begin{theorem}[D.2.2; Proved existence and exact Euler number]\label{st:D.2.2}
The class of compact complex threefolds with $b_1=b_2=0$,
a proper connected-fibre map with smooth two-torus general fibre,
and trivial rational monodromy is nonempty. Every member satisfies
\[
                   c_3(Y)=0,\qquad b_3(Y)=2.
\]
Every individual reduced fibre support has Euler number zero and
the geometric description of C.2.G.1. Arbitrary singularities and
full-fibre multiplicities are allowed. There is an example with
integral $H_1(Y,\mathbb Z)=H_2(Y,\mathbb Z)=0$ and
algebraic dimension one.
\end{theorem}
\begin{proof}
The unrestricted lower bound of C.2.G.1 and the independent
trivial-monodromy upper bound of D.2.5 give $c_3(Y)=0$.
Poincar\'e duality gives $b_3(Y)=2-c_3(Y)=2$.
C.2.G.1 expresses $c_3$ as the sum of the nonnegative
Euler numbers of the reduced fibre supports, so each summand
vanishes and satisfies that theorem's exact equality criterion.
Statement D.2.7 constructs the asserted example with integral
homology vanishing and algebraic dimension one. D.2.3--4 also
give a separate local proof of the equality when all supports
have normal crossings.
\end{proof}

\begin{proposition}[D.2.3; Proved logarithmic period obstruction]
\label{st:D.2.3}
Let $h:U\to\Delta$ be a proper holomorphic map from a smooth
complex threefold, with connected fibres, smooth two-torus fibres
over $\Delta^*$, and central fibre $F=\sum_i m_iD_i$ with
normal-crossings reduced support $D=\bigcup_iD_i$.
If every normalization component of $D$ has geometric genus zero,
the monodromy on $H^2(U_t;\C)$ is nontrivial. In particular the
monodromy on $H^1(U_t;\Q)$ is nontrivial. No K\"ahler assumption
on $U$ or $D$ is made, and the positive integers $m_i$ are arbitrary.
\end{proposition}
\begin{proof}
We first identify exactly the analytic comparison used in the
argument. Put
\[
\mathscr K^\bullet=
 \Omega_U^\bullet(\log D)\big/
 \bigl(h^*\Omega_\Delta^1(\log0)\wedge
       \Omega_U^{\bullet-1}(\log D)\bigr),\qquad
\mathscr H^2=\mathbb R^2h_*\mathscr K^\bullet.
\]
In \cite[\S2.1, Theorem~2.18, \S2.19, Proposition~2.20 and
Theorem~2.21]{Steenbrink1976}, the assumptions for these assertions
are precisely a proper holomorphic map from a connected complex
manifold to a disc, smooth off zero, with normal-crossings support.
They precede the additional Hodge-theoretic assumptions of the
later sections. They give a locally free bundle $\mathscr H^2$,
its actual Gauss--Manin connection
\[
\nabla:\mathscr H^2\longrightarrow
\Omega_\Delta^1(\log0)\otimes\mathscr H^2,
\]
and the comparison
$\mathscr H^2|_{\Delta^*}\simeq
R^2h_*\C\otimes\OO_{\Delta^*}$.
The residue eigenvalues belong to $[0,1)$ and the extended
monodromy at zero is
$\exp(-2\pi i\operatorname{Res}_0\nabla)$.
Its canonical base-change map identifies the fibre at zero with
$\mathbb H^2(D,\mathscr K^\bullet/t\mathscr K^\bullet)$.
The local calculation used below is the same monomial Koszul
calculation as \cite[Proposition~1.13 and proof of
Proposition~2.20]{Steenbrink1976}; we give it explicitly for the
subcomplex needed here. No local invariant-cycle surjectivity
or polarization is invoked.

Suppose that the monodromy on $H^2$ were the identity. Its
residue has all eigenvalues zero, and its exponential is the
identity. A nilpotent matrix whose exponential is the identity
is zero. Hence the residue vanishes. In a holomorphic frame
$\nabla=d+A(t)\,dt/t$ with $A(0)=0$, so this is a holomorphic
connection across zero. It has a holomorphic horizontal frame
on the disc. Consequently the coordinates of every holomorphic
section of $\mathscr H^2$ in a flat frame, that is its periods,
are holomorphic and bounded near zero.

The determinant of the locally free relative logarithmic complex is
\[
\mathscr K^2\simeq
 \omega_U(D)\otimes h^*\omega_\Delta(0)^{-1}
 \simeq\omega_{U/\Delta}(D-F).
\]
The direct image $h_*\mathscr K^2$ is a torsion-free
coherent sheaf of rank one, because a smooth torus fibre has
one-dimensional space of holomorphic two-forms. Shrink the disc
and choose a generator $\Omega$ of this line bundle.
Thus $\Omega$ determines a holomorphic section $[\Omega]$ of
$\mathscr H^2$ by the inclusion of the last term of the
relative de Rham complex. Its smooth-fibre comparison is the
ordinary class $[\Omega_t]\in H^2(U_t;\C)$.

First suppose that $\Omega$ restricts nontrivially to a normalized
component $T$ of $D$. There is a canonical isomorphism
\[
 \mathscr K^1|_T\simeq\Omega_T^1(\log D_T),\qquad
 \mathscr K^2|_T\simeq\omega_T(D_T),
\]
where $D_T$ is the reduced conductor. To specify the map, on the
branch $z_i=0$ subtract
$\operatorname{res}_i(\eta)(dt/t)/m_i$ from an absolute logarithmic
one-form $\eta$ and restrict its tangential part. The normal
residue of $dt/t$ is $m_i$, so this kills exactly the subbundle
divided out in $\mathscr K^1$ and is independent of the lift.
The exterior square gives the second isomorphism. If every
residue of $\Omega|_T$ vanished, its local logarithmic expression
would be holomorphic on $T$. This contradicts $p_g(T)=0$.
Thus there is a smooth point of a double curve, away from triple
points, where its residue is nonzero.

In a small chart about that point, write $h=x^ay^b=t$ and use $z$
along the double curve. The relative two-form is
\[
 \Omega_t=g(x,y,z)\,\frac{dx}{x}\wedge dz,
 \qquad g(0,0,z_0)\ne0.
\]
Choose a fixed small $z$-disc and $\epsilon>0$ such that
$|g(x,y,z)|\ge c>0$ whenever $|x|,|y|<\epsilon$ there.
The map to the annulus
$|t|^{1/a}\epsilon^{-b/a}<|x|<\epsilon$
has $b$ sheets. Integration on its inverse image and the fixed
$z$-disc gives
\[
 \int_{U_t}\Omega_t\wedge\overline{\Omega_t}
 \ \ge\ c'\log\frac{\epsilon^{a+b}}{|t|}
 \ \longrightarrow\ +\infty.
\]
Here the integrand is a nonnegative real volume form on the
complex surface; its local expression is a positive constant
times $|g|^2|x|^{-2}\,dA(x)dA(z)$.

On the other hand, $\Omega_t$ is closed, and this integral is
the fixed cup-product pairing of $[\Omega_t]$ with its complex
conjugate. All its periods are bounded by the holomorphic
connection argument, so the pairing is bounded. This is a
contradiction.

It remains to treat the case where $\Omega$ vanishes on every
branch of $D$. Put
\[
 \mathscr J=\OO_U(-D),\qquad
 \mathscr B^\bullet=
 \mathscr J\mathscr K^\bullet/t\mathscr K^\bullet.
\]
Both $\mathscr J\Omega_U^\bullet(\log D)$ and
$t\Omega_U^\bullet(\log D)$ are subcomplexes for the absolute
exterior derivative: locally the logarithmic derivatives of
$z_1\cdots z_r$ and $t=\prod z_i^{m_i}$ are holomorphic
linear combinations of the $d\log z_i$. Hence the locally split
absolute/relative logarithmic exact sequence remains an exact
sequence of complexes after tensoring by $\mathscr J/(t)$:
\[
0\longrightarrow\mathscr B^\bullet[-1]
 \xrightarrow{\ (dt/t)\wedge\ }
 (\mathscr J/(t))\Omega_U^\bullet(\log D)
 \longrightarrow\mathscr B^\bullet\longrightarrow0.
\tag{D.18}
\]
Signs in the shifted differential are chosen so that the first
arrow is a morphism of complexes. Its connecting morphism
$\rho_{\mathscr B}:\mathscr B^\bullet\to\mathscr B^\bullet$
in the derived category defines degree-preserving endomorphisms
on both cohomology sheaves and hypercohomology. The inclusion
$\mathscr J/(t)\hookrightarrow\OO_U/(t)$ gives a commutative
diagram from (D.18) to the analogous exact sequence for
$\mathscr K^\bullet/t\mathscr K^\bullet$. The induced
endomorphism for the latter is precisely the Gauss--Manin
residue. Thus the map on hypercohomology intertwines these
endomorphisms.

We compute the eigenvalues of $\rho_{\mathscr B}$ without
assuming reduced multiplicities. At a point of $r$ branches,
write
\[
 t=\prod_{i=1}^r z_i^{m_i},\qquad
 e=\gcd(m_1,\ldots,m_r),\qquad
 y=\prod_{i=1}^r z_i^{m_i/e},\qquad t=y^e.
\]
The transverse holomorphic Poincar\'e lemma reduces the relative
complex to the Koszul complex with commuting operators
\[
 D_i=z_i\partial_{z_i}
           -\frac{m_i}{m_1}z_1\partial_{z_1},
 \qquad 2\le i\le r,
\]
on the convergent-power-series module
$(z_1\cdots z_r)/(z_1^{m_1}\cdots z_r^{m_r})$.
The monomials diagonalize these operators. A monomial with some
nonzero eigenvalue has a contracting Koszul homotopy, obtained
by division by that eigenvalue. These divisions preserve
convergence: every nonzero eigenvalue has absolute value at
least $1/m_1$. The surviving monomials have exponent vector
$a(m_1/e,\ldots,m_r/e)$. Membership in $\mathscr J$ requires
$a\ge1$, and quotienting by $t$ requires $a<e$. Consequently
each cohomology stalk is spanned by
\[
 y^a\,d\log z_{i_1}\wedge\cdots\wedge d\log z_{i_q},
 \qquad 1\le a<e,
\]
subject to the relative relation $\sum m_i d\log z_i=0$.
Since $dy^a=(a/e)y^a\,dt/t$, the connecting endomorphism
acts on these representatives by $a/e$. In particular zero is
not an eigenvalue; when $e=1$ the complex is acyclic.

The hypercohomology spectral sequence
\[
 H^p\bigl(D,\mathcal H^q(\mathscr B^\bullet)\bigr)
 \ \Longrightarrow\ \mathbb H^{p+q}(D,\mathscr B^\bullet)
\]
is equivariant for this derived endomorphism. There are only
finitely many branch multiplicities, so all eigenvalues on its
terms and abutment lie in a finite subset of $(0,1)$. For
example the product of the corresponding linear factors
annihilates every cohomology sheaf, and a power annihilates the
abutment. Hence $\rho_{\mathscr B}$ is invertible on
$\mathbb H^2(D,\mathscr B^\bullet)$. The residue on
$\mathscr H^2|_0$ is zero. Equivariance therefore proves that
\[
 \mathbb H^2(D,\mathscr B^\bullet)\longrightarrow
 \mathbb H^2(D,\mathscr K^\bullet/t\mathscr K^\bullet)
       =\mathscr H^2|_0
 \quad\hbox{is the zero map}.
\tag{D.19}
\]

The central class of $\Omega$ comes from the left side of
(D.19), because $\Omega\in\mathscr J\mathscr K^2$.
Thus $[\Omega]$ vanishes at zero in the locally free bundle
$\mathscr H^2$, and $[\Omega]/t$ is a holomorphic section of
that bundle. Its periods are bounded. On a smooth fibre this
section is the de Rham class of $\Omega_t/t$.

Primitivity of the generator $\Omega$ implies that its order
$b_i$ along at least one $D_i$ is less than $m_i$; otherwise
$\Omega/t$ would be a holomorphic section of $\mathscr K^2$
over $U$, contradicting its choice as a generator of
$h_*\mathscr K^2$. Choose a smooth point of that component
away from all other branches and all further zeros. There
are coordinates with $t=z^{m_i}$ in which
\[
 \Omega=z^{b_i}a(z,\xi,\eta)\,d\xi\wedge d\eta,
 \qquad a(0,0,0)\ne0.
\]
On one sheet above a fixed small $(\xi,\eta)$-polydisc, the
positive integral of $\Omega_t/t$ against its conjugate has
the lower bound
\[
 \int_{U_t}\frac{\Omega_t}{t}\wedge
                    \overline{\frac{\Omega_t}{t}}
 \ \ge c''|t|^{2(b_i/m_i-1)}\longrightarrow+\infty.
\]
Its expression as the fixed cup-product pairing of bounded
periods again gives a contradiction. This proves nontriviality
of the $H^2$ monodromy in the remaining case.

Finally the natural
cup-product isomorphism
$\bigwedge^2H^1(U_t;\Q)\simeq H^2(U_t;\Q)$ for a two-torus
is monodromy-equivariant, proving the final assertion.
\end{proof}

\begin{theorem}[D.2.4; Proved sign in the full normal-crossings class]
\label{st:D.2.4}
For the threefolds and maps in D.2 whose reduced fibre supports
all have normal crossings, the Euler number is
\[
c_3(Y)=0,\qquad b_3(Y)=2.
\]
All multiplicities, including unequal multiplicities on a
reducible support, are allowed. In this class there are no
positive-Euler-number examples with trivial monodromy.
\end{theorem}
\begin{proof}
Statement C.2.NC.5 gives $c_3(Y)=2r$, where $r$ counts the
irreducible fibre supports with rational normalization. A fibre
of that kind is necessarily $mS$ for one positive integer $m$.
Choose a small disc around its value and perform the normalized
base change $t=u^m$. Statement C.2.NC.3 identifies the resulting
map to the original neighbourhood as finite \emph{\'etale} of
degree $m$, and its central fibre $S'$ as reduced with normal
crossings. The total space is therefore smooth. Its nearby
fibres are the original connected tori. Stein factorization
over the normal disc shows that its central fibre is connected
as well: the finite Stein factor has generic degree one and is
the disc itself.

Normalization commutes with this \emph{\'etale} base change.
Each normalized component of $S'$ is an unramified finite cover
of the smooth rational normalization $T$ of $S$. Since $T$ is
simply connected, these normalized components are copies of
$T$, and all have geometric genus zero. The local monodromy
after base change is the $m$th power of the original monodromy,
so remains the identity on $H^1$. This contradicts D.2.3.
Therefore $r=0$, proving $c_3(Y)=0$. Poincar\'e duality and
$b_1(Y)=b_2(Y)=0$ give $c_3(Y)=2-b_3(Y)$, hence $b_3(Y)=2$.
The example of D.2.7 realizes this normal-crossings class.
\end{proof}
